\section{Radiometric Camera Model}
\label{sec:radiometric-camera-model}
\label{sec:image-formation}
Let $\mathcal{I}\subseteq\mathbb{R}^2$ be the sensor plane. A pixel
$\vbu$ with the area $\mathcal{I}_{\vbu}\subseteq\mathcal{I}$ measures
the radiant flux~\cite{pharr2023physically}
\begin{equation}
\label{eq:per-channel-rad-flux}
    \Phi(\vbu) \coloneqq \int_{\mathcal{I}_{\vbu}}\int_{\calOmega(\vbp)} L(\vbx(\vbp,\vbomega), \vbomega_o(\vbp,\vbomega)) \cos\theta(\vbomega)\odif{\vbomega}\odif{\vbp},
\end{equation}
where $\calOmega(\vbp)$ is the solid angle under which the sensor
point $\vbp$ sees the aperture, $\vbx(\vbp,\vbomega)$ the scene point
visible from $\vbp$ along the direction $\vbomega$,
$L(\vbx,\vbomega_o)$ the radiance that leaves $\vbx$ in the direction
$\vbomega_o$ towards the lens, and $\theta$ the angle of incidence on
the sensor. The inner integral is the irradiance $E(\vbp)$. All
quantities are per color channel, and we omit the channel index
(\cref{sec:supp-image-formation} gives the spectral formulation). The
radiance $L$ follows from the lighting and the reflectance of the
scene (\cref{eq:scene-rad}). The lens model determines
$\calOmega(\vbp)$, $\vbx(\vbp,\vbomega)$ and $\vbomega_o(\vbp,\vbomega)$. With a pinhole aperture,
$\calOmega(\vbp)$ is a single direction, so every scene point is in
focus and there is no vignetting (\cref{sec:supp-pinhole-radiometry}).

\subsection{Projection}
Let $\vbpi^{-1}: \mathcal{I}\rightarrow\mathcal{S}^2$ be the
back-projection of a central geometric camera model with the center
of projection $\vbo$. It maps a sensor point $\vbp$ to the unit
direction $\hat{\vbd}(\vbp)\coloneqq\vbpi^{-1}(\vbp)$ of its ray. For
a model that is given by its forward projection, we invert the
projection numerically per sample. The radiometric camera keeps
$\vbpi^{-1}$ as it is.

\subsection{Aperture and Focus Surface}
\label{sec:defocus}
\begin{figure}[t]
        \centering
        % fig_model.tex — the lift of a central geometric model, for a perspective and for a fisheye projection (hand-written TikZ).
% The geometric model sets the direction d(p) at the center of projection o; the lift adds the aperture, whose rays meet in
% the point of focus x_f(p) on the focus surface F.
\usetikzlibrary{intersections}%
\resizebox{\linewidth}{!}{%
\begin{tikzpicture}[x=1cm, y=1cm, >=stealth, font=\footnotesize,
    ray/.style={line width=0.45pt, blue!70!black},
    chief/.style={line width=0.75pt, black},
    inside/.style={line width=0.3pt, densely dashed, black!55},
    axis/.style={line width=0.3pt, densely dotted, black!60},
    dim/.style={line width=0.3pt, black!60}]
  % ---------- the parts both panels share
  \newcommand{\rccamera}{%
    \draw[axis] (-0.25,0) -- (1.75,0);
    \draw[line width=1pt] (0,-1.0) -- (0,1.0);
    \node[anchor=south] at (0,1.0) {$\mathcal{I}$};
    \draw[line width=1.8pt] (1.0,0.45) -- (1.0,1.0) (1.0,-0.45) -- (1.0,-1.0);
    \node[anchor=south] at (1.0,1.0) {$\mathcal{A}$};
    \coordinate (p) at (0,0.6); \coordinate (o) at (1.0,0);
    \coordinate (a1) at (1.0,0.45); \coordinate (a2) at (1.0,-0.45);
    \draw[inside] (p) -- (a1) (p) -- (a2) (p) -- (o);
    \fill (p) circle (1.2pt) node[anchor=east, inner sep=2pt] {$\vbp$};
    \fill (o) circle (1.2pt);
    \fill[blue!70!black] (a2) circle (1.2pt);}
  % ---------- left: a perspective projection; the ray leaves o along the line from p
  \begin{scope}
    \rccamera
    \node[anchor=north east, inner sep=1.5pt] at (o) {$\vbo$};
    \node[anchor=north east, inner sep=1.5pt] at (a2) {$\vba$};
    \coordinate (xf) at (2.75,-1.05);
    \draw[line width=0.9pt, red!75!black] plot[smooth, tension=0.7] coordinates {(2.62,1.25) (2.78,0.1) (2.75,-1.05) (2.60,-1.85)};
    \node[red!75!black, anchor=south] at (2.62,1.25) {$\mathcal{F}$};
    \path[name path=sceneA] plot[smooth, tension=0.8] coordinates {(3.95,1.25) (3.75,0.4) (3.92,-0.5) (3.74,-1.3) (3.9,-2.05)};
    \path[overlay, name path=rA0] (o) -- ($(o)!2.2!(xf)$);
    \path[overlay, name path=rA1] (a1) -- ($(a1)!2.2!(xf)$);
    \path[overlay, name path=rA2] (a2) -- ($(a2)!2.2!(xf)$);
    \path[name intersections={of=sceneA and rA0, by=s0}];
    \path[name intersections={of=sceneA and rA1, by=s1}];
    \path[name intersections={of=sceneA and rA2, by=s2}];
    \draw[chief] (o) -- (s0);
    \draw[chief, ->] (o) -- ($(o)!0.42!(xf)$);
    \draw[dim] (1.92,0.12) -- ($(o)!0.30!(xf)$);
    \node[anchor=south west, inner sep=1pt] at (1.80,0.10) {$\hat{\vbd}(\vbp)$};
    \draw[ray] (a1) -- (s1) (a2) -- (s2);
    \draw[ray, ->] (a2) -- ($(a2)!0.5!(xf)$);
    \node[blue!70!black, anchor=north, inner sep=2pt] at (1.75,-0.82) {$\vbpi^{-1}(\vbp,\vba)$};
    \draw[line width=1pt] plot[smooth, tension=0.8] coordinates {(3.95,1.25) (3.75,0.4) (3.92,-0.5) (3.74,-1.3) (3.9,-2.05)};
    \node[anchor=south] at (3.9,1.25) {scene};
    \draw[line width=2.2pt, blue!70!black, line cap=round] (s1) -- (s0) -- (s2);
    \draw[dim] (3.30,-2.02) -- ($(s0)+(-0.07,-0.05)$);
    \node[blue!70!black, anchor=north, inner sep=1.5pt] at (3.05,-1.98) {region seen by $\vbp$};
    \fill[red!75!black] (xf) circle (1.7pt);
    \node[anchor=south west, inner sep=1pt] at ($(xf)+(0.10,0.16)$) {$\vbx_f(\vbp)$};
    \node[anchor=north] at (1.9,-2.45) {lifted perspective};
  \end{scope}
  % ---------- right: a fisheye projection; the model turns the ray at o, and the points of focus surround the camera
  \begin{scope}[xshift=4.75cm]
    \rccamera
    \draw[axis] (o) -- ++(1.15,-0.69);  % the continuation of the line from p: the direction of a perspective model
    \coordinate (xf) at ($(o)+(-62:1.45)$);
    \draw[line width=0.9pt, red!75!black] ($(o)+(62:1.45)$) arc[start angle=62, end angle=-92, radius=1.45];
    \node[red!75!black, anchor=south west, inner sep=1pt] at ($(o)+(58:1.45)$) {$\mathcal{F}$};
    \path[name path=sceneB] plot[smooth, tension=0.8] coordinates {($(o)+(42:2.25)$) ($(o)+(12:2.1)$) ($(o)+(-18:2.25)$) ($(o)+(-48:2.1)$) ($(o)+(-76:2.25)$) ($(o)+(-100:2.15)$)};
    \path[overlay, name path=rB0] (o) -- ($(o)!2.4!(xf)$);
    \path[overlay, name path=rB1] (a1) -- ($(a1)!2.4!(xf)$);
    \path[overlay, name path=rB2] (a2) -- ($(a2)!2.4!(xf)$);
    \path[name intersections={of=sceneB and rB0, by=t0}];
    \path[name intersections={of=sceneB and rB1, by=t1}];
    \path[name intersections={of=sceneB and rB2, by=t2}];
    \draw[chief] (o) -- (t0);
    \draw[chief, ->] (o) -- ($(o)!0.5!(xf)$);
    \draw[dim] (1.72,0.12) -- ($(o)!0.32!(xf)$);
    \node[anchor=south west, inner sep=1pt] at (1.60,0.10) {$\hat{\vbd}(\vbp)$};
    \draw[ray] (a1) -- (t1) (a2) -- (t2);
    \draw[line width=1pt] plot[smooth, tension=0.8] coordinates {($(o)+(42:2.25)$) ($(o)+(12:2.1)$) ($(o)+(-18:2.25)$) ($(o)+(-48:2.1)$) ($(o)+(-76:2.25)$) ($(o)+(-100:2.15)$)};
    \draw[line width=2.2pt, blue!70!black, line cap=round] (t1) -- (t0) -- (t2);
    \fill[red!75!black] (xf) circle (1.7pt);
    \node[anchor=west, inner sep=3pt] at (xf) {$\vbx_f(\vbp)$};
    \node[anchor=north] at (1.9,-2.45) {lifted fisheye};
    \node[anchor=south west, inner sep=1pt] at ($(o)+(42:2.25)$) {scene};
  \end{scope}
\end{tikzpicture}}%

        \caption{The radiometric camera lifts any central geometric model, here a perspective and a fisheye model, in a cross-section. The geometric model alone sets the direction $\hat{\vbd}(\vbp)$ in which the ray of the sensor point $\vbp$ leaves the center $\vbo$ of the aperture $\mathcal{A}$; the model does not trace the light inside the camera (dashed). The lift adds the aperture: the rays through all aperture points $\vba$ meet in the point of focus $\vbx_f(\vbp)$ on the focus surface $\mathcal{F}$, which is close to a plane for the perspective model and surrounds the camera for the fisheye. Where the scene is off $\mathcal{F}$, the cone of these rays meets it in a region (thick blue), from which $\vbp$ collects its light.}
        \label{fig:model}
\end{figure}
The thin lens~\cite{pharr2023physically} replaces the pinhole by an
aperture disc $\mathcal{A}$ of radius $r$, centered at $\vbo$ and
parallel to the sensor. All rays from a sensor point $\vbp$ through
$\mathcal{A}$ meet in one point $\vbx_f(\vbp)$ on the plane of focus.
Where the scene does not pass through $\vbx_f(\vbp)$, the cone of
these rays intersects it in a region and not in a point, and $\vbp$
collects the light of this whole region, which defocuses the image. For a real
lens, the points of focus form a surface $\mathcal{F}$
(\cref{fig:model}) that is not a plane in general \cite{Hecht+2023}: it is approximately a plane for a
thin-lens-like lens and a sphere for a fisheye
\cite{Kingslake2010,wide-field-of-view-2017}. We therefore learn the
focus surface. We represent it by the inverse distance $\tau(\vbp)$
from $\vbo$ to the point of focus $\vbx_f(\vbp)\in\mathcal{F}$ along
the ray of $\vbp$,
\begin{equation}
    \vbx_f(\vbp) \coloneqq \vbo + \frac{1}{\tau(\vbp)}\,\hat{\vbd}(\vbp),
\end{equation}
and extend the back-projection to the aperture
(\cref{fig:model}). The ray from $\vbp$
through an aperture point $\vba\in\mathcal{A}$ is directed from
$\vba$ towards $\vbx_f(\vbp)$:
\begin{equation}
\label{eq:lifted-back-projection}
    \vbpi^{-1}(\vbp,\vba)
    \coloneqq \frac{\vbx_f(\vbp) - \vba}{\lVert \vbx_f(\vbp) - \vba \rVert}
    = \frac{\tau(\vbp)(\vbo-\vba) + \hat{\vbd}(\vbp)}{\lVert\tau(\vbp)(\vbo-\vba) + \hat{\vbd}(\vbp)\rVert}.
\end{equation}
This is the lift of the geometric model. Setting the aperture $\mathcal{A}=\{\vbo\}$ reduces the lifted model back to the geometric one, since $\vbpi^{-1}(\vbp,\vbo)=\vbpi^{-1}(\vbp)$. Where the latter back-projects to a single ray,
$\vbpi^{-1}(\vbp,\mathcal{A})$ back-projects to a cone of rays
through $\vbx_f(\vbp)$. Similarly, with the perspective back-projection and a plane of focus, the model reduces to
the thin lens (\cref{sec:supp-thin-lens-equiv}). The inverse distance
avoids intersecting the rays with $\mathcal{F}$, and unlike the
distance itself, it keeps the gradient from vanishing for a far focus
\cite{unified-inverse-depth,direct-sparse-odometry}
(\cref{sec:supp-inverse-depth}). We model $\tau$ by a thin plate
spline \cite{thin-plate-spline} over a grid restricted to the illuminated region
of the sensor.

\subsection{Vignetting}
\label{sec:vignetting}
Written as an integral over the aperture, the flux of the thin lens
carries the factor $\cos^4\theta/s'^2$ that expresses its vignetting, with the image distance $s'$~\cite{Goossens_2019}
(\cref{sec:supp-thin-lens-vignetting}). Real lenses differ much
from it (\cref{fig:vignetting}). Hence, we replace this factor by a generic
vignetting function $V$. We also normalize by the aperture area $\pi r^2$,
so that the aperture radius changes the defocus but not the
brightness of the image, such that its scale is carried by the lighting (see \cref{sec:lighting}, \ref{sec:supp-gauge-freedoms}) alone. To spare us computing a boundary integral (see \cref{sec:boundary-terms}), we substitute the aperture point
$\vba=\vbo+r\vbxi$ with $\vbxi$ in the unit disc $\mathcal{D}$ of the
aperture plane. The area element $\odif{\vba}=r^2\odif{\vbxi}$
cancels the $r^2$ of the normalization, and the radiometric camera
measures
\begin{equation}
\label{eq:rad-flux-with-vignetting}
    \Phi(\vbu)
    = \frac{1}{\pi}\int_{\mathcal{I}_{\vbu}}V(\vbp)
    \int_{\mathcal{D}}L\bigl(\vbx(\vbp,\vba),-\vbpi^{-1}(\vbp,\vba)\bigr)\odif{\vbxi}\odif{\vbp},
\end{equation}
where $\vbx(\vbp,\vba)$ is the scene point that the ray from $\vba$ in
the direction $\vbpi^{-1}(\vbp,\vba)$ meets. The radius now
appears only in the position of the aperture point, and no longer in
the domain of the integral. We assume a vignetting
that is radially symmetric about the principal point
$\vbc\in\mathcal{I}$ and represent it by a cubic B-spline
\cite{piegl1997nurbs}
\begin{equation}
    V(\vbp) \coloneqq \sum_{i=0}^{n} N_{i,3}\Bigl(\frac{\lVert \vbp-\vbc \rVert}{r_{\text{max}} + \lVert \vbp-\vbc \rVert}\Bigr)P_{i},
\end{equation}
with the basis functions $N_{i,3}$, the control points
$\vbTheta_v = \{P_{i}\}_{i\in\{0,\dots,n\}}$, $n=15$, and the largest
distance $r_{\text{max}}$ from the principal point to an illuminated
pixel. Its knots are uniform in the image
(\cref{sec:supp-vignetting-uniform-knots}).

% one-panel caption (fig_model_single.tex): \caption{The radiometric camera in a cross-section. The geometric model gives the ray $\hat{\vbd}(\vbp)$ of the sensor point $\vbp$ through the center $\vbo$ of the aperture $\mathcal{A}$ of radius $r$. The rays from $\vbp$ through all aperture points $\vba$ meet in the point of focus $\vbx_f(\vbp)$ on the focus surface $\mathcal{F}$, at the distance $1/\tau(\vbp)$ from $\vbo$. Where the scene is off $\mathcal{F}$, the cone of these rays intersects it in a region (thick blue), from which $\vbp$ collects its light.}
